% !TEX root = cv-en.tex
% !TEX program = xelatex
% English CV content. Edit this file for wording, dates, and entries.
\def\cvpdftitle{hy30nq | Security Research CV}
\def\cvfooterlabel{hy30nq · Security Research CV}
\def\cvupdated{Updated September 2026}
% Shared monochrome resume layout for cv-ko.tex and cv-en.tex.
% Edit typography, spacing, and entry macros in this file.
\documentclass[10pt,a4paper]{article}

\usepackage[a4paper,left=18mm,right=18mm,top=15mm,bottom=15mm]{geometry}
\usepackage{fontspec}
\usepackage{array}
\usepackage{tabularx}
\usepackage{enumitem}
\usepackage{fancyhdr}
\usepackage{hyperref}
\usepackage{xcolor}
\usepackage{ragged2e}
\usepackage{needspace}

% Local macOS preference, followed by portable fallbacks for XeLaTeX.
\IfFontExistsTF{NanumSquare Neo}{
  \setmainfont{NanumSquare Neo}[BoldFont={NanumSquare Neo ExtraBold}]
  \setsansfont{NanumSquare Neo}[BoldFont={NanumSquare Neo ExtraBold}]
}{
  \IfFontExistsTF{Pretendard}{
    \setmainfont{Pretendard}
    \setsansfont{Pretendard}
  }{
    \IfFontExistsTF{Noto Sans CJK KR}{
      \setmainfont{Noto Sans CJK KR}
      \setsansfont{Noto Sans CJK KR}
    }{
      \setmainfont{Arial Unicode MS}
      \setsansfont{Arial Unicode MS}
    }
  }
}

% Strict black-and-white palette.
\definecolor{ink}{HTML}{000000}

\hypersetup{
  unicode=true,
  hidelinks,
  pdftitle={\cvpdftitle},
  pdfauthor={hy30nq},
  pdfsubject={Security research, vulnerability analysis, and CTF engineering}
}

\pagestyle{fancy}
\fancyhf{}
\renewcommand{\headrulewidth}{0pt}
\renewcommand{\footrulewidth}{0pt}
\fancyfoot[C]{%
  \makebox[\textwidth]{%
    \fontsize{7.2}{9}\selectfont
    \cvfooterlabel\hfill\cvupdated\ · \thepage
  }%
}
\setlength{\footskip}{8mm}

\setlength{\parindent}{0pt}
\setlength{\parskip}{0pt}
\setlength{\tabcolsep}{0pt}
\renewcommand{\arraystretch}{1.0}
\hyphenpenalty=10000
\exhyphenpenalty=10000
\color{ink}

\setlist[itemize]{
  leftmargin=4.2mm,
  itemsep=1.6pt,
  topsep=2.3pt,
  parsep=0pt,
  partopsep=0pt,
  label=\raisebox{0.15ex}{\tiny\textbullet}
}

\newcommand{\cvheader}[1]{%
  \begin{tabularx}{\textwidth}{@{}X>{\RaggedLeft\arraybackslash}p{66mm}@{}}
    {\fontsize{30}{32}\selectfont\bfseries hy30nq} &
    \begin{tabular}[t]{@{}r@{}}
      \fontsize{8.15}{10.6}\selectfont hy30nq@gmail.com\\
      \fontsize{8.15}{10.6}\selectfont github.com/hy30nq03\\
      \fontsize{8.15}{10.6}\selectfont hyeonql.tistory.com
    \end{tabular}\\[-0.2mm]
    {\fontsize{11.2}{13.6}\selectfont\bfseries #1} &
    {\fontsize{8.15}{10.6}\selectfont hy30nq.com}\\
  \end{tabularx}
  \vspace{2.6mm}
  \hrule height 0.8pt
}

\newcommand{\sectiontitle}[2]{%
  \needspace{18mm}
  \vspace{5.2mm}%
  \begin{tabularx}{\textwidth}{@{}Xr@{}}
    {\fontsize{12.8}{15}\selectfont\bfseries #1} &
    {\fontsize{7.2}{9}\selectfont\bfseries\MakeUppercase{#2}}
  \end{tabularx}
  \vspace{1.1mm}
  \hrule height 0.45pt
  \vspace{2.8mm}
}

\newcommand{\entry}[4]{%
  \needspace{14mm}
  \begin{tabularx}{\textwidth}{@{}>{\RaggedRight\arraybackslash}p{24mm}>{\RaggedRight\arraybackslash}X@{}}
    {\fontsize{8.1}{10.2}\selectfont #1} &
    {\fontsize{9.8}{12.2}\selectfont\bfseries #2}\\[-0.2mm]
    & {\fontsize{8.0}{10.2}\selectfont\itshape #3}\\[0.55mm]
    & {\fontsize{8.65}{11.3}\selectfont #4}
  \end{tabularx}
  \vspace{3.1mm}
}

\newcommand{\compactentry}[3]{%
  \needspace{6mm}
  \begin{tabularx}{\textwidth}{@{}>{\RaggedRight\arraybackslash}p{24mm}>{\RaggedRight\arraybackslash}X@{}}
    {\fontsize{8.1}{10.2}\selectfont #1} &
    {\fontsize{8.75}{11.1}\selectfont #2\enspace
     {\fontsize{7.9}{9.8}\selectfont · #3}}
  \end{tabularx}
  \vspace{1.7mm}
}

% Compact bullet list for CTF authoring and participation history.
\newenvironment{cvlist}{%
  \begin{itemize}[
    leftmargin=4.2mm,
    itemsep=2.0pt,
    topsep=0.5pt,
    parsep=0pt,
    partopsep=0pt
  ]
  \fontsize{8.75}{11.1}\selectfont
}{%
  \end{itemize}
}

\newcommand{\tag}[1]{%
  \begingroup
  \setlength{\fboxsep}{1.2mm}%
  \setlength{\fboxrule}{0.4pt}%
  \fbox{\strut\fontsize{7.7}{9.5}\selectfont\bfseries #1}%
  \endgroup
}

\newcommand{\cvskills}{%
  \vfill
  \begin{center}
    \tag{Vulnerability Research}\hspace{1.4mm}
    \tag{CTF Engineering}\hspace{1.4mm}
    \tag{Security Education}\hspace{1.4mm}
    \tag{Python · Docker · Git}
  \end{center}
}


\begin{document}
\sffamily

\cvheader{Security Researcher · CTF Builder}

% =============================================================================
% Affiliations
% =============================================================================
\sectiontitle{Affiliations}{Education \& Communities}

\compactentry{2023--Present}
  {\textbf{Korea University · B.S. in AI Cybersecurity}}
  {Privacy \& Security Convergence · Integrated B.S.-M.S. Program}

\compactentry{Present}
  {\textbf{Science \& Technology Specialist Officer Program Candidate}}
  {Cohort 3}

\compactentry{}
  {\textbf{SeeCure Lab · KUality · Knights of the SPACE · Rubiyalab}}
  {Security Research \& CTF Communities}

% =============================================================================
% Awards: hyperlinks are used only in this section
% =============================================================================
\sectiontitle{Awards}{Recognition}

\compactentry{2025}
  {\href{https://security.korea.ac.kr/m06_04/59}{\textbf{Encouragement Prize} · 2nd Privacy Protection \& Security Policy Idea Contest}}
  {Korea University School of Cybersecurity}

\compactentry{2024}
  {\href{https://blog.naver.com/kisiaedu/223785868393}{\textbf{Excellence Award · KISIA Chairman's Award} · AI Security Program Essay Contest}}
  {Korea Information Security Industry Association}

\compactentry{2024}
  {\href{https://www.kisia.or.kr/talent_support/education_business/ai_security_technology_development/}{\textbf{Minister of Science and ICT Award · Top Graduate} · KISIA AI Security Program}}
  {Top overall evaluation across coursework, labs, and project work}

\compactentry{2024}
  {\href{https://www.boannews.com/media/view.asp?idx=132213\&direct=mobile}{\textbf{Grand Prize (1st) · Minister of Science and ICT Award} · KISIA Security Developer Hackathon}}
  {Chrome extension for malicious-URL detection}

\compactentry{2024}
  {\href{https://www.msit.go.kr/bbs/view.do?sCode=user\&mId=113\&mPid=238\&bbsSeqNo=94\&nttSeqNo=3184227}{\textbf{Top 20 · Minister of Science and ICT Certificate} · WhiteHat School Cohort 1}}
  {Selected among the top 20 of 320 participants}

% =============================================================================
% CTF finals
% =============================================================================
\sectiontitle{CTF Finals}{Competition Results}

\begin{cvlist}
  \item \textbf{2026 · DEF CON 34 CTF} · Finalist · Team SSS
  \item \textbf{2025 · WhiteHat Contest} · Finalist · 5th in the qualifier · Team KUality
  \item \textbf{2025 · Black Hat MEA CTF} · Finalist · 41st in the qualifier · Team PPU-AENG Overflow
\end{cvlist}

% =============================================================================
% CTF challenge authoring: add one bullet per event, newest first
% =============================================================================
\sectiontitle{CTF Challenge Authoring}{Operations \& Review}

\begin{cvlist}
  \item \textbf{2026 · KCTF 2026 (KUality)} · Challenge authoring · Liberation Signal 1--3
  \item \textbf{2026 · CSKWS Joint Seminar CTF} · Event operations and challenge authoring · Ghost Display, Black Cache
  \item \textbf{2026 · HACKTHEON Sejong Student Cybersecurity Festival} · CTF operations and challenge authoring
  \item \textbf{Aug 2026 · Information Security Fundamentals Program} · Introductory instruction and Mini CTF operations
  \item \textbf{Feb 2026 · INC0GNITO FESTIVAL CTF} · Event operations and qualifier/final challenge authoring · Project Incoherence
  \item \textbf{Feb 2026 · 5th Hanyang Hacking Defense Contest, HCTF with HSPACE} · Challenge authoring
  \item \textbf{Aug 2025 · SSD CTF} · Challenge authoring
\end{cvlist}

\newpage

% =============================================================================
% Projects
% =============================================================================
\sectiontitle{Projects}{Security Engineering}

\compactentry{2026--Present}
  {\textbf{Adversarial Attacks and Incident Traceability for AI Systems}}
  {Participating Researcher · Ongoing}

\compactentry{2026--Present}
  {\textbf{Security Analysis of Remote Device-Management Protocols (TR-069) II}}
  {Participating Researcher · Ongoing}

\compactentry{2026--Present}
  {\textbf{AI-Based Future-Threat Scenario Prediction from Threat-Actor Campaigns and MITRE D3FEND-Based Countermeasure Enhancement}}
  {Participating Researcher · Ongoing}

\entry{2026}
  {Tool/MCP Call Gateway for Securing AI IDEs}
  {Project Semester · Vulnerability Research}
  {Analyzed and reproduced unsafe shell-command execution from untrusted project Hook configurations in Cursor IDE without workspace-trust enforcement or sandboxing.
   Prepared a reproducible PoC and responsible-disclosure materials, then designed a Verify--Restrict--Log gateway methodology for mitigation.}

\entry{Jun--Jul 2025}
  {HSPACE Penetration-Testing Curriculum}
  {Content Developer · Knights of the SPACE}
  {Built a mock e-commerce service and tiered penetration-testing scenarios so beginners could learn complete attack flows in a repeatable lab.}

\entry{Jun--Aug 2025}
  {CTF Platform Based on Real Bug-Bounty Cases}
  {Backend Developer · FinderGap}
  {Implemented the backend of a training platform that converts real bug-bounty cases into reproducible CTF challenges.}

\entry{2024}
  {Automated WordPress Plugin Vulnerability Analysis}
  {Korea University AI Cybersecurity Symposium}
  {Co-developed an automated vulnerability detection tool for WordPress plugins and implemented its SQL injection analysis module.
   Presented the methodology and implementation at the departmental symposium.}

\entry{2024}
  {AI-Powered Malicious URL Classifier}
  {Frontend · Data Preprocessing · KISIA AI Security Network}
  {Contributed to phishing-URL data cleaning and the model-training pipeline, and developed the frontend for a Chrome extension that detects malicious URLs.}

% =============================================================================
% CVE
% =============================================================================
\sectiontitle{CVE}{Vulnerability Disclosure}

\entry{2024}
  {CVE-2024-11950 · XnSoft XnView Classic RWZ Integer Underflow}
  {Independent Vulnerability Research · ZDI-24-1640 · High Severity}
  {Discovered and responsibly disclosed an integer underflow and resulting memory corruption in XnView Classic's RWZ parser.
   The issue was published by ZDI and assigned CVE-2024-11950 with remote-code-execution impact.}

% =============================================================================
% Other
% =============================================================================
\sectiontitle{Other}{Additional Experience}

\entry{Mar--Aug 2025}
  {Teaching Assistant, KITRI WhiteHat School Cohort 3}
  {Stage 1 · Midterm · Stage 3}
  {Supported hands-on labs and midterm presentations, providing participant feedback and program coordination.}

\entry{Jun 2024--Present}
  {KnockOn Q\&A Mentor}
  {Security Community Mentor}
  {Guide learners through web-security and vulnerability-analysis questions and hands-on practice.}

\compactentry{2024}
  {\textbf{Information Processing Craftsman}}
  {National Technical Qualification}

\compactentry{}
  {\textbf{Languages} · Korean (native) · English (reading/writing) · Chinese (HSK 4, HSKK Elementary)}
  {Languages}

\cvskills

\end{document}
